\section{二审程序}
\begin{enumerate}
  \item 两审终审。
  \item 原告、被告、有独三，担责的无独三可以上诉。\al{检察院二审抗诉 \cuo}。
  \item 上诉对象 \yjt \zx{非一审终审、未生效的一审判决，3 类裁定}。\bs{不包括调解书 \yjt 一审终审}！！
  \item 上诉期 \yjt \zx{判决 15 日，裁定 10 日}。
  \item \bs{必须提交书面上诉状，口头无效}。给谁都行，但一审法院\zd{没有审查权}（只是\zd{跑腿}的。\al{递情书}）。
        \al{区法院发现甲的上诉已超过上诉期 \yjt 上诉状和上诉材料一并移送上级法院}。
  \item 当事人的诉讼地位。
        \begin{enumerate}
          \item \zx{谁上诉谁就是上诉人}，对谁提谁是被上诉人。\al{甲乙丙均不服判决，提出上诉。都是上诉人，笑死}。
          \item 无独三能做为被上诉人。
          \item 把一审判决划出来，假设上诉人的上诉请求完全成立，二审法院直接照上诉请求改，
                \zx{对谁不利谁就是被上诉人，其他人按照原审地位列明}。\al{一审判道歉，二审判不用道歉}。
        \end{enumerate}
  \item 程序。由审判员组成合议庭，没有陪审员。可以适用独任制的条件见基本制度。
  \item 以开庭审理为原则 \yjt 特殊情况可以不开庭。\al{阅卷、调查、询问当事人、没有提出新的事实理由、法院认为}。
  \item 审理范围。
        \begin{enumerate}
          \item 基于处分原则，\bs{仅限上诉请求范围}！！除非违背公序良俗。\al{判了 AB，上诉 A，即使 B 有错，也只审 A}。
          \item 基于两审终审，\bs{仅限于一审审理过的问题}（判了啥）！！\al{请求 AB，判了 A，上诉请求 A，新增 C}。
                \begin{enumerate}
                  \item 遗漏的请求/必须参加的当事人、一审不给离，二审给离。\bs{只准调不准判}，调解不成，发回重审。
                        \begin{enumerate}
                          \item \al{一审判决遗漏一个继承人}。
                          \item \al{不得对\xhx{遗漏的请求}进行审理 \cuo{调解也是审理的一种方式}}。
                          \item \al{只能对\xhx{上诉请求}进行\xhx{审理并作出判决} \dui{微妙的区别}}。
                        \end{enumerate}
                  \item 增加/变更诉讼请求、反诉。调解不成，告知另诉。
                  \item 小技巧。一审法院或多或少\zx{有错的/没错的} \yjt 调解不成 \yjt \zx{发回重审/告知另诉}。
                  \item \bs{仅告知另诉的情形，经当事人同意}，\zx{可以一起审理，并作出判决}。\al{BC 当事人同意，可一并裁判 \cuo}。
                  \item \zd{既有情形 1，又有情形 2 \yjt 分别处理}。
                \end{enumerate}
        \end{enumerate}
  \item 结案有 4 种方式。\zx{裁判、调解、撤回起诉、撤回上诉}。
        \begin{enumerate}
          \item 裁判。
                \begin{enumerate}
                  \item 原判判决/裁决\zd{正确} \yjt 判决/裁定\zd{驳回}上诉，\zd{维持}原判决/裁定。\bs{文书要一一对应}。
                  \item 事实 \hz 法律错误 \yjt 判决/裁定改判。\zd{有错就改}。二审法院心里有答案了。
                  \item 基本事实\zd{不清} \yjt 要查清事实 \yjt \zd{可改可发}。
                  \item 严重违反法定\zd{程序} \yjt 只能撤销原判，发回重审（两审终审制）。\zx{逆向考，选项说应当发回重审，其实是在问你程序有没有错}。
                  \item \bs{发回重审，一次为限（没有例外）}。民事诉讼的终极目的是要解决纠纷化解矛盾的。
                \end{enumerate}
          \item 调解。\bs{必须}制作调解书，\bs{送达后原判决视为撤销}。\al{应当写明撤销原判、申请恢复对原判决的执行 \cuo}。
          \item \zx{撤回起诉}（地基。我不告了） vs \zx{撤回上诉}（二楼。我服一审判决）。
                \begin{enumerate}
                  \item \zx{双方同意 + 法院审查}是否损害国社三利益。
                        \al{发现和解协议内容与原判\xhx{认定事实不一致}，不允许撤诉 \cuo{不属于国社三情况}}。
                  \item 裁定准予\zx{撤回上诉之日起，一审判决生效。只能向一审申请执行一审判决}。\al{申请执行和解协议 \cuo}。
                  \item 裁定\zx{撤回起诉时，应当一并裁定撤销原一审判决}（从头撤到尾）。处分起诉权。
                  \item 二审、再审\zx{撤回起诉/上诉不允许再次起诉} \yjt 已经实体处理。
                        \al{甲起诉乙违约，一审判决乙败诉，乙提起二审，双方达成和解，约定一起撤诉（甲 \yjt 撤回起诉，乙 \yjt 撤回上诉）}。
                  \item 没交上诉费，视为撤回上诉。
                \end{enumerate}
        \end{enumerate}
\end{enumerate}
